% 05-scaling — Scaling to the full-parallel geometry and a predictive ASIC.
% Every number carries a `% src:` comment naming the docs/ report and heading it was read from.
% Nothing here is a new measurement; nothing predictive is called "the chip".

\section{Scaling to the full-parallel geometry and a predictive ASIC}
\label{sec:scaling}

The shipped 16/48 core (\S\ref{sec:silicon}) decodes a worst-case window in 2085
cycles, \SI{15.64}{\micro\second} at the KV260's \SI{133.3}{\mega\hertz}.
% src: docs/perf/q7-02-ngroup-sweep.md "What this measures and why"
Latency is cycles divided by clock, so the two questions this section answers
are how far the cycle count can be pushed down by widening the core, and how
far the clock can be pushed up by leaving the FPGA. Both were measured rather
than modelled, and two of the three answers came out worse than the plan
assumed. We report them in the order they were obtained.

\subsection*{The banking family is the only implementable one}

Two architectures span the design space between the banked core and a fully
unrolled one. \texttt{bp\_relay\_unroll\_pipe} stamps $1/\mathit{NGROUP}$ of
the check and variable slots and time-multiplexes each across $\mathit{NGROUP}$
groups per BP phase. Its cycle count is exactly affine,
$\mathrm{cycles}(\mathit{NGROUP}) = 122\cdot\mathit{NGROUP}+240$, and the
co-simulation is 40/40 bit-identical to the golden at every setting measured.
% src: docs/perf/q7-02-ngroup-sweep.md "The cycle half: closed form, bit-exact everywhere"
The knob nevertheless has no feasible setting. Five values were synthesised
out-of-context on the KV260 part at a \SI{5.0}{\nano\second} target: every one
is over the device budget \emph{and} slower than the banked core. The cheapest
member, $\mathit{NGROUP}=144$ with one check slot and six variable slots
stamped, still needs 703\,698 LUTs (601\,\% of the device) at
\SI{17.5}{\mega\hertz}; the fastest, $\mathit{NGROUP}=16$, reaches
\SI{132.8}{\micro\second} at 1474\,\% of the device, 8.5$\times$ slower than
the shipped core. F$_{\max}$ is flat at 16.5--\SI{17.5}{\mega\hertz} across a
9$\times$ span of the knob.
% src: docs/perf/q7-02-ngroup-sweep.md "The area/Fmax half: measured"
The reason is structural: each slot gathers across $\mathit{NGROUP}$ groups and
there are $\texttt{BP\_C}/\mathit{NGROUP}$ slots, so the mux volume is
$\mathit{NGROUP}$-invariant by construction and the crossbar is simultaneously
the area floor and the critical path. Least squares gives
$\mathrm{LUT} = 615\,985 + 18\,415\,474/\mathit{NGROUP}$ within 5.7\,\%.
% src: docs/perf/q7-02-ngroup-sweep.md "Why: the crossbar is NGROUP-invariant by construction"
Banking, which replaces the flat register array plus crossbar with
index-addressed banks, is therefore not an optimisation of a working parallel
design but the thing that makes it implementable.
% src: docs/perf/q7-02-ngroup-sweep.md "Consequences"

Within the banked family the plan had assumed cycles scale linearly with bank
count. Cycle-measuring five geometries showed that they do not. The counts fit
the exact closed form
\begin{equation}
  \mathrm{cycles} = 60\,(G_C+G_V) + 420 + (2G_V + G_C),
  \label{eq:cycles}
\end{equation}
where $G_C,G_V$ are the check and variable group counts and the constant
$420 = 60\cdot 7$ is a per-iteration pipeline tail (seven drain cycles in each
of the $\mathrm{LEGS}\cdot\mathrm{ITERS}=60$ iterations) that exists regardless
of banking. Only the first term shrinks with banking, so 4$\times$ banking
(16/48 $\to$ 64/192) yields 2.28$\times$ fewer cycles, 2085 $\to$ 913, and the
tail grows from 20\,\% to 47\,\% of the budget.
% src: docs/qec/asic-architecture.md "Banking-scaling qualification (measured, Q7-08 follow-up)"
The floor is the full-parallel geometry, $G_C=G_V=1$: 543 cycles. This
configuration could not be generated when the qualification was written; three
generator defects (routing-network lanes sized by bank count rather than tap
count, a zero-width row address at $\lceil\log_2 1\rceil=0$, and ROM literals
beyond Verilator's 65\,536-bit limit) were fixed in Task~B0 Option~A, after
which 48/288, 144/432 and 144/864 all pass the 40-shot golden 40/40 at 789, 605
and 543 cycles, with the header for every previously qualified geometry
byte-identical apart from two guard lines.
% src: docs/qec/open-silicon-program.md "Task B0 Option A RESULT (2026-07-30)"
Since the 420-cycle tail is $\mathrm{LEGS}\cdot\mathrm{ITERS}\cdot 7$, the
dominant remaining latency lever once banking is exhausted is the iteration
budget, ITERS, which is set by the decoder's LER at the operating point rather
than by the hardware.
% src: docs/qec/asic-architecture.md "Banking-scaling qualification (measured, Q7-08 follow-up)"

\subsection*{A large FPGA: it fits, and the clock collapses}

Both 64/192 and 144/864 were implemented end to end (synthesis, opt, place,
phys\_opt, route) on the AMD Virtex UltraScale+ HBM \texttt{xcvu47p-fsvh2892-2-e}
of AWS's F2 instances, with Vivado 2025.2 on a \texttt{z1d.2xlarge} build
instance, at a \SI{5.0}{\nano\second} target, default directives, no
floorplanning and no fanout constraints. The instance lived 7.3 hours and cost
about \$5.50; 64/192 implemented in 65 minutes, 144/864 in 5\,h\,40\,min.
% src: docs/perf/q7-02-fullparallel-fpga.md "8. What was run"
No FPGA instance was rented: the build only needs the Vivado Enterprise licence
that the FPGA Developer AMI bundles, and a bitstream is never loaded.
% src: docs/qec/b2-aws-build-runbook.md "What you are renting and why"
Table~\ref{tab:fpga} gives the post-route result.

\begin{table}[t]
  \centering
  \footnotesize
  \caption{Post-route implementation of two banked geometries on the VU47P
    (Vivado 2025.2, \SI{5.0}{\nano\second} target, default directives). The
    shipped KV260 row is the silicon measurement of \S\ref{sec:silicon}.}
  \label{tab:fpga}
  \begin{tabular}{@{}lrrrrr@{}}
    \toprule
    Geometry & Cycles & LUTs & Util. & F$_{\max}$ & Latency \\
             &        &      & (\%)  & (MHz)      & (\si{\micro\second}) \\
    \midrule
    64/192   & 913  & 247\,434 & 19.0 & 150.4 & 6.07 \\
    144/864  & 543  & 994\,700 & 76.3 & 97.3  & 5.58 \\
    \addlinespace
    16/48, KV260 & 2085 & --- & --- & 133.3 & 15.64 \\
    \bottomrule
  \end{tabular}
  % src: docs/perf/q7-02-fullparallel-fpga.md "9. Results"
\end{table}

The fit gate passes: 76.3\,\% is under the 90\,\% criterion, BRAM and LUTRAM
are at zero and registers at 9.4\,\%, so LUTs are the only resource that
matters. The latency case for full-parallel on FPGA does not pass:
1.68$\times$ fewer cycles at a 1.55$\times$ slower clock is a net 1.09$\times$
for 4.0$\times$ the area. 64/192 delivers 92\,\% of the latency in a quarter of
the device.
% src: docs/perf/q7-02-fullparallel-fpga.md "9. Results"
The clock loss is not logic depth. The worst path has eight logic levels and a
\SI{10.251}{\nano\second} data-path delay of which \SI{0.750}{\nano\second}
(7.3\,\%) is logic and \SI{9.501}{\nano\second} (92.7\,\%) is routing: a
control counter, \texttt{pc\_reg[1]}, reaching a min-sum cell across a die large
enough to need three super logic regions, with Laguna SLR-crossing tiles inside
the congested windows. It is the fanout-7941 control net the out-of-context
probe had flagged, and every standard mitigation (\texttt{MAX\_FANOUT}, per-bank
replication, SLR-aware floorplanning) is untried, so \SI{97.3}{\mega\hertz} is
a floor for this design rather than its ceiling.
% src: docs/perf/q7-02-fullparallel-fpga.md "10. Why the clock collapsed"
Two calibration lessons are worth carrying forward. The same RTL synthesised to
803\,518 CLB LUTs on the KV260 part with Vivado 2024.2 and to 1\,000\,035 on the
VU47P with 2025.2, 24\,\% more, so cross-part LUT extrapolation is worth about
$\pm25\,\%$; and post-route utilisation moved less than 1\,\% off
post-synthesis on this part.
% src: docs/perf/q7-02-fullparallel-fpga.md "11. Where Step 0's projections were wrong"
% src: docs/perf/q7-02-fullparallel-fpga.md "9. Results"
Sub-microsecond needs \SI{543}{\mega\hertz} on 543 cycles; nothing in this
family on any FPGA is within a factor of three of that.
% src: docs/perf/q7-02-fullparallel-fpga.md "4. Calibration"

\subsection*{7\,nm predictive: the geometry penalty is mostly an FPGA artefact}

The FPGA result raised a specific risk: the only ASIC-node clock this project
had measured, \SI{686.13}{\mega\hertz}, came from the 16/48 geometry, and the
program had been quoting 543 cycles against it as if the clock were
geometry-independent.
% src: docs/perf/q7-02-fullparallel-fpga.md "12. What this means for the ASIC"
Both geometries were therefore run through OpenROAD-flow-scripts
\citep{openroad} on the ASAP7 predictive 7\,nm platform \citep{clark2016asap7},
with the Q7-08 latch register file, identical knobs
(\texttt{CORE\_UTILIZATION}=45, \texttt{PLACE\_DENSITY}=0.55, SDC period
\SI{1000}{\pico\second}, timing repair skipped because \texttt{repair\_timing}
segfaults on this netlist class), and RC-extracted post-route timing. The
144/864 run took about 20 hours on an EPYC box and completed with $rc=0$.
% src: docs/perf/q7-02-b3-asap7-fullparallel.md "1. What was run"
Table~\ref{tab:asap7} summarises both.

\begin{table}[t]
  \centering
  \small
  \caption{ASAP7 7\,nm predictive, post-route, same flow and knobs for both
    geometries. F$_{\max}$ is \texttt{report\_clock\_min\_period}, setup-only,
    on a netlist whose hold violations are unrepaired (\S\ref{sec:limitations}).
    Power is from a gate-level VCD of the 16/48 netlist at the
    \SI{1}{\giga\hertz} SDC; it was not measured at 144/864.}
  \label{tab:asap7}
  \begin{tabular}{@{}lrr@{}}
    \toprule
    & 16/48 & 144/864 \\
    \midrule
    Cycles (worst case)                & 2085   & 543 \\
    Design area at 50\,\% util. (\si{\micro\metre\squared}) & 81\,090 & 434\,450 \\
    Die (\si{\milli\metre\squared})    & 0.163  & 0.869 \\
    F$_{\max}$, setup-only (MHz)       & 686.13 & 614.59 \\
    WNS at \SI{1000}{\pico\second} (ps) & $-622.73$ & $-1117.97$ \\
    Latency (\si{\micro\second})       & 3.04   & 0.88 \\
    DRC violations                     & 0      & 0 \\
    Hold violations (unrepaired)       & 44\,704 & 35\,616 \\
    \addlinespace
    Power, steady BP (W)               & 0.149  & --- \\
    Power, idle, clock running (W)     & 0.138  & --- \\
    Energy per 2085-cycle window (\si{\micro\joule}) & 0.31 & --- \\
    \bottomrule
  \end{tabular}
  % src: docs/perf/q7-02-b3-asap7-fullparallel.md "2. Results, against the 16/48 baseline"
  % src: docs/perf/q7-02-asap7-timing.md "7. Closure note" (0.163 mm^2 die)
  % src: docs/perf/q7-02-b3-asap7-fullparallel.md "3.2 Neither run had timing repair" (hold counts)
  % src: docs/perf/q7-02-asap7-timing.md "4. Power" (power rows)
\end{table}

The clock is geometry-dependent on the ASIC too, but only mildly: the step from
the small to the full-parallel geometry costs $-10.4\,\%$ on ASAP7 against
$-35.3\,\%$ on the VU47P. An ASIC has neither fixed routing tracks nor SLR
boundaries, and its buffer trees absorb the high-fanout control net that
dominated the FPGA path; roughly a third of the penalty transferred. Area scales
sub-linearly, 5.36$\times$ for 8.5$\times$ the logic, where the FPGA paid
8.53$\times$ in CLB LUTs; of the placer's $+11.18\,\%$ over the raw netlist,
routability inflation was $+1.28\,\%$ and timing-driven growth $+9.90\,\%$.
% src: docs/perf/q7-02-b3-asap7-fullparallel.md "2. Results" ("The geometry penalty", "Area scales sub-linearly")
The headline, $543\ \text{cycles} / \SI{614.59}{\mega\hertz} =
\SI{0.88}{\micro\second}$, is a sub-microsecond worst-case window on a 7\,nm
predictive node, DRC-clean, at \SI{0.869}{\milli\metre\squared}.
% src: docs/perf/q7-02-b3-asap7-fullparallel.md "4. Verdict"
Figure~\ref{fig:latency} places it against the FPGA points.

We state the two qualifications here rather than only in
\S\ref{sec:limitations}. First, every F$_{\max}$ in Table~\ref{tab:asap7} is a
\emph{setup-only} figure: hold was never repaired, the 16/48 netlist carries
44\,704 hold violations (worst \SI{-745.9}{\pico\second}, 98\,\% on the D pins
of the latch register file) and the 144/864 netlist 35\,616 (worst
\SI{-1486.9}{\pico\second}). Hold does not scale with the clock period, so
these netlists would not decode correctly at any frequency; the first
gate-level co-simulation of the 16/48 netlist fails 40/40 in exactly the way
STA predicts, and \SI{686}{\mega\hertz} is the speed of the datapath, not of a
working part. A separately repaired but unrouted run gave
\SI{527.66}{\mega\hertz}, so the defensible 16/48 band is
528--\SI{686}{\mega\hertz}.
% src: docs/perf/q7-02-asap7-timing.md "1. Headline: the Fmax number straddles the budget"
% src: docs/perf/q7-02-asap7-timing.md "4b. Gate-level co-simulation"
Second, ASAP7 is a proxy: the margin below \SI{1}{\micro\second} is 12\,\%, a
1.2$\times$ node penalty gives \SI{1.06}{\micro\second}, and nothing in this
work measures the gap to a production node.
% src: docs/qec/open-silicon-program.md "4. Risk register" ("The node")

\subsection*{Power with real switching activity}

The default-activity power ORFS reports, \SI{0.656}{\watt}, is not a decoder
number. The routed 16/48 netlist (639\,322 cells, among them 55\,703 latches
and 43\,025 flops) was compiled with Verilator against the ASAP7 behavioural
cell models and driven by the same testbench and 40 golden vectors the RTL
passes; two 300-cycle VCD windows, one inside a steady-state BP iteration and
one with the decoder parked, were fed to OpenSTA with the extracted SPEF,
annotating 2\,128\,987 pin activities with none left unannotated.
% src: docs/perf/q7-02-asap7-timing.md "4. Power — measured with real switching activity (Task A4)"
At the \SI{1}{\giga\hertz} SDC the core draws \SI{0.149}{\watt} during BP and
\SI{0.138}{\watt} idle; the default activity was 4.4$\times$ too high, almost
entirely in the combinational group (\SI{0.466}{\watt} $\to$
\SI{0.0115}{\watt}). Ninety-three per cent of the real power is clock tree
plus sequential internal: the ASAP7 platform marks integrated clock-gating
cells \texttt{DONT\_USE}, so every register toggles its clock pin every cycle
and the datapath itself is about \SI{10}{\milli\watt}. Energy per worst-case
window is \SI{149}{\pico\joule} per cycle $\times$ 2085 cycles $=$
\SI{0.31}{\micro\joule}, clock-independent to first order and inside the
\SI{1}{\micro\joule} budget with 3$\times$ margin, at 7\,nm predictive; the
activity comes from the not-quite-bit-exact simulation of the previous
paragraph, and absolute watts on a predictive NLDM library should be read as
$\pm30\,\%$.
% src: docs/perf/q7-02-asap7-timing.md "4. Power — measured with real switching activity (Task A4)"

\begin{figure}[t]
  \centering
  \includegraphics[width=\linewidth]{latency.pdf}
  \caption{Worst-case decode latency against banking geometry on the three
    platforms measured: KV260 silicon (16/48), VU47P post-route (64/192,
    144/864) and ASAP7 7\,nm predictive post-route (16/48, 144/864). ASAP7
    points are setup-only clocks on hold-unrepaired netlists. Data as in
    Tables~\ref{tab:fpga} and~\ref{tab:asap7}.}
  \label{fig:latency}
\end{figure}
